\subsection{Analytic inversion of the complete contact action}
\label{subsec:v59-analytic-inverse}

The diagonal estimate does not control the powers in
\eqref{eq:v58-triangular-factorization}.  For analytic contact germs a
stronger conclusion follows from a different representation of the
\emph{complete} derivative.  Stationarity removes the variation of the
orbit.  The initial contact then contributes an invertible multiplication
operator, while every later contact evaluates the boundary variation at a
strictly contracted argument.  This gives an inverse in an analytic
function norm, including all lower-degree couplings.  The norm, the local
nature of the result, and the restriction on physical collars are specified
below; none of the smooth or statistical norms used earlier is silently
replaced by it.

Write $D_R=\{z\in\mathbb C:|z|<R\}$ and give
$\mathscr B_R=H^\infty(D_R;\mathbb C^2)$ the norm
$\|f\|_R=\max_b\sup_{D_R}|f_b|$.  For $m\ge3$ let
\begin{equation}\label{eq:v59-analytic-spaces}
 \mathscr X_{R,m}=\{f\in\mathscr B_R:f_b^{(k)}(0)=0
                  \text{ for }0\le k<m,\ b=0,1\},
 \qquad \mathscr X_R=\mathscr X_{R,3}.
\end{equation}
These are closed complex Banach spaces.  Real germs mean the fixed points
of $f(z)\mapsto\overline{f(\bar z)}$.  Fix $g>0$ and a real analytic
anchored pair $\psi^*=(\psi_0^*,\psi_1^*)$, holomorphic near
$\overline{D_{R_0}}$, with $(\psi_b^*)''(0)=\kappa_b>0$.
The affine space $\psi^*+\mathscr X_R$ fixes the gap and the two quadratic
jets.  Denote its half-line action map by
\[
                 \mathscr S(\psi)=(S_{0,\psi},S_{1,\psi}).
\]
Here the action is the actual stationary sum
\eqref{eq:v4-boundary-action}, not a formally assigned series.

\begin{lemma}[Holomorphic half-lines on a fixed function space]
\label{lem:v59-holomorphic-halfline}
There are $0<R<R_0$, an open ball $\mathscr U$ about $\psi^*$ in
$\psi^*+\mathscr X_R$, and $0<\rho<1$ for which both half-line orbits
are holomorphic in $(\psi,u)$, with
\begin{equation}\label{eq:v59-disc-contraction}
 x_{b,0}(u)=u,\qquad |x_{b,i}(u)|\le\rho^i|u|\quad(i\ge1,
                         \ u\in D_R,\ \psi\in\mathscr U).
\end{equation}
The stationary sums define a holomorphic map
$\mathscr U\to\mathscr S(\psi^*)+\mathscr X_R$.  For real members
of $\mathscr U$ these are the physical local actions on a common smaller
real collar.  Positivity of curvature is required on that smaller collar,
not on the whole complex disc or at its boundary.
\end{lemma}

\begin{proof}
We provide the function-space details because pointwise analytic dependence
of a finite family would not suffice.  Put $t=e^{-\gamma}$ and
$\lambda_b=\mathfrak r_bt$.  By
\eqref{eq:v58-admissible-contraction},
$\lambda_*:=\max(\lambda_0,\lambda_1)<1$.  The linear orbit is
$x^{\rm lin}_{b,i}(u)=\ell_{b,i}u$, where
$\ell_{b,2k}=t^{2k}$ and
$\ell_{b,2k+1}=\lambda_bt^{2k}$; hence
$|\ell_{b,i}|\le\lambda_*^i$.  Choose
\[
                   \lambda_*<\beta<\rho<\sigma<1.
\]
On sequences of bounded holomorphic functions vanishing at zero use
\[
 \|y\|_\beta=\sup_{i\ge1}\beta^{-i}\|y_i\|_{H^\infty(D_R)}.
\]
The fixed-reference Green kernel \eqref{eq:v4-half-green} is bounded
in this norm.  Indeed its absolute entries are bounded by a constant
times $t^{|i-k|}$, and
$\sup_i\sum_{k\ge1}t^{|i-k|}\beta^{k-i}<\infty$ since $t<\beta<1$.

Consider $\psi=\psi^*+\eta$ with
$\eta\in\mathscr X_R$ and $\|\eta\|_R<\varepsilon R^3$.
The maximum principle applied to $\eta_b(z)/z^3$ gives
$|\eta_b(z)|\le\varepsilon|z|^3$ on $D_R$.
On $D_{\sigma R}$, Cauchy's estimates give uniform third-derivative
bounds, and, for $k\le2$,
\begin{equation}\label{eq:v59-protected-cauchy}
 |\eta_b^{(k)}(z)|\le C_{\sigma,k}\varepsilon|z|^{3-k}.
\end{equation}
For example, write $\eta_b=z^3v_b$, use
$\|v_b\|_R\le\varepsilon$, and differentiate the product using
Cauchy bounds for $v_b$ on a slightly larger inner disc.  The reference
pair has the same type of remainder bounds with a fixed constant.

The stationarity equation at an interior site of type $b_i$ is
\begin{equation}\label{eq:v59-interior-equation}
 \frac{h_{i-1}\psi_{b_i}'(x_i)+x_i-x_{i-1}}{\ell_{i-1}}
 +\frac{h_i\psi_{b_i}'(x_i)+x_i-x_{i+1}}{\ell_i}=0,
\end{equation}
where $h_i=g+\psi_{b_i}(x_i)+\psi_{1-b_i}(x_{i+1})$ and
$\ell_i=\sqrt{h_i^2+(x_{i+1}-x_i)^2}$, with the square-root branch
near $g$.  Crucially, when $x_0=u$ approaches the boundary of $D_R$,
this equation uses only $\psi_b(u)$, never $\psi_b'(u)$.
Every differentiated graph is evaluated at an interior site, whose
argument lies in $D_{\sigma R}$.  Thus an unbounded differentiation
operator on $H^\infty(D_R)$ is not being used.

Subtract the linear Jacobi equation, whose diagonal is $2c_{b_i}/g$
and adjacent entries are $-1/g$.  On the weighted ball
$\|x-x^{\rm lin}\|_\beta\le C_0R^2$, the nonlinear residual has
norm at most $C R^2$ and Lipschitz constant at most $CR$ in the
interior sequence.  These estimates follow by expanding
\eqref{eq:v59-interior-equation}, using
\eqref{eq:v59-protected-cauchy} at interior arguments and the cubic
value remainder at $u$.  Locality bounds neighboring weights by fixed
factors $\beta^{\pm1}$; products of two such coordinates have norm
$O(R^2)$.  All denominators stay in a fixed neighborhood of $g$.
Choose $C_0$ larger than twice the residual constant times the Green
norm, then choose $R$ small.  The Green fixed-point map preserves this
ball and has contraction constant less than $1/2$.

This construction is consistent with the protected inner arguments:
for every $i\ge1$, the correction vanishes at zero, so the Schwarz
bound gives
\[
 |x_{b,i}(u)|\le
       (\lambda_*^i+C_0R\beta^i)|u|\le\rho^i|u|
\]
after $R$ is decreased.  The last inequality follows uniformly in $i$
from $\lambda_*/\rho+C_0R\beta/\rho<1$.
Compositions and graph derivatives at the protected arguments are
holomorphic maps of the Banach variables: their Cauchy expansions
converge uniformly while these arguments remain inside $D_{\sigma R}$.
The fixed-point iterates therefore converge locally uniformly as
Banach-valued holomorphic maps.  Their limit is holomorphic.  Shrink the
parameter ball once to keep these bounds and its inverse margin uniform.

Each action summand is bounded by $CR^2\rho^{2i}$, so their sum is a
locally uniformly convergent holomorphic map into $\mathscr B_R$.
Its constant, linear and quadratic jets are respectively zero, zero and
the fixed $a_b$ of \eqref{eq:v4-boundary-hessian}.  Action differences
belong to $\mathscr X_R$.  On $|u|<R/2$, Cauchy bounds imply that real
nearby graphs remain strictly convex and separated.  The real
stationary sequence just constructed is the unique small half-line of
Lemma~\ref{lem:v4-halfline}, hence gives its actual local action there.
No assertion that every perturbed germ extends to an entire periodic
obstacle is required for this local statement.
\end{proof}

\begin{lemma}[The complete derivative as weighted boundary visits]
\label{lem:v59-envelope-operator}
On a smaller ball as in Lemma~\ref{lem:v59-holomorphic-halfline}, put
\[
 w_{b,i}(u)=\frac{g+\psi_{b+i}(x_{b,i}(u))+
                           \psi_{b+i+1}(x_{b,i+1}(u))}
                         {\ell_{b,i}(u)},\qquad
 \alpha_b=w_{b,0},\qquad v_{b,i}=w_{b,i-1}+w_{b,i}\ (i\ge1),
\]
where types are read modulo two.  There are constants $W,a_*>0$ such
that $|w_{b,i}|\le W$ and $|\alpha_b|\ge a_*$ on $D_R$.
For every $\eta\in\mathscr X_R$, the full Fr\'echet derivative is
\begin{equation}\label{eq:v59-full-envelope}
 (L_\psi\eta)_b(u):=(D\mathscr S(\psi)\eta)_b(u)
 =\alpha_b(u)\eta_b(u)
       +\sum_{i\ge1}v_{b,i}(u)\eta_{b+i}(x_{b,i}(u)).
\end{equation}
Write $L_\psi=A_\psi+K_\psi$ for the multiplier and the sum.
Both preserve every $\mathscr X_{R,m}$, and
\begin{equation}\label{eq:v59-tail-bound}
 \|A_\psi^{-1}\|\le a_*^{-1},\qquad
 \|K_\psi|_{\mathscr X_{R,m}}\|
                    \le\frac{2W\rho^m}{1-\rho^m}.
\end{equation}
The operator $K_\psi$ is compact on $\mathscr X_R$; specifically it
has a rank at most $2(m-3)$ approximation with error at most
$2W\rho^m/((1-\rho)(1-\rho^m))$.
\end{lemma}

\begin{proof}
All length factors are close to $g$ and $w_{b,i}$ is close to one on
the chosen complex domain, so the stated bounds hold.  Differentiate a
finite action sum along $\psi+s\eta$.  The internal orbit derivatives
cancel by \eqref{eq:v59-interior-equation}; the derivative of $x_0$ is
zero.  The only remaining orbit term is
\[
 \partial_2\ell_{b+N-1}(x_{b,N-1},x_{b,N})\,
                           D_\psi x_{b,N}[\eta].
\]
On a smaller parameter ball, the holomorphic fixed-point construction
gives $\|D_\psi x_{b,N}[\eta]\|_R\le C_R\beta^N\|\eta\|_R$.
The first factor is at most $CR\rho^{N-1}$.
The terminal term therefore tends to zero in $H^\infty(D_R)$.
The direct variation of each flight is exactly
$w_{b,i}\{\eta_{b+i}(x_{b,i})+
\eta_{b+i+1}(x_{b,i+1})\}$.  These variations are norm summable:
$\eta\in\mathscr X_R$ and the Schwarz estimate give a majorant
$C\|\eta\|_R\rho^{3i}$.  Passing to the limit in the finite
identity proves \eqref{eq:v59-full-envelope}.  The initial visit occurs
once and all internal visits occur twice, as in the finite-jet proof.

For $\eta\in\mathscr X_{R,m}$ the same Schwarz estimate is
$|\eta_b(z)|\le\|\eta\|_R(|z|/R)^m$.  Summing it at the
contracted arguments gives \eqref{eq:v59-tail-bound}.  Nonvanishing
of $\alpha_b$ makes its reciprocal a bounded holomorphic multiplier.
It preserves the order of vanishing.

For compactness, in every interior evaluation replace $\eta_b$ by its
Taylor polynomial through degree $m-1$.  Each of the $2(m-3)$ scalar
coefficients then multiplies one fixed pair of holomorphic functions,
obtained by summing $v_{b,i}x_{b,i}^n$ over the appropriate parity.
Those sums converge in norm.  The resulting operator has the stated
finite rank.  By Cauchy's coefficient bound, the evaluation remainder
at $x_{b,i}$ is at most
$\|\eta\|_R\rho^{im}/(1-\rho^i)$.
Summation proves the asserted error, which tends to zero.  Notice that
no continuation of the weights or orbits beyond $D_R$ is used.
Compactness here is approximation by finite-rank maps on the same disc;
it is not a gain of holomorphic radius for the full derivative.
\end{proof}

\begin{theorem}[Local analytic inversion of the full action map]
\label{thm:v59-analytic-contact-inverse}
For every real analytic positive-curvature anchored pair $\psi^*$ and
gap $g>0$ as above, one can choose $R>0$ and neighborhoods for which
\[
 \mathscr S:\psi^*+\mathscr X_R
                 \longrightarrow\mathscr S(\psi^*)+\mathscr X_R
\]
is a biholomorphism between those neighborhoods.  In particular, on a
sufficiently small convex ball there are constants $0<c<C<\infty$ with
\begin{equation}\label{eq:v59-nonlinear-bilip}
 c\|\psi-\widetilde\psi\|_R
 \le\|\mathscr S(\psi)-\mathscr S(\widetilde\psi)\|_R
 \le C\|\psi-\widetilde\psi\|_R.
\end{equation}
These constants control the complete map, not only pairs with fixed lower
jets.  For every $N\ge3$, its derivative and inverse induce bounded
mutually inverse operators on
$\mathscr X_R/\mathscr X_{R,N+1}$, with bounds independent of $N$
in the quotient norms induced by $\|\cdot\|_R$.

For $0<r<R$, all contact coefficients of two such pairs satisfy
\begin{equation}\label{eq:v59-all-coefficients}
 \sum_{n\ge3}\frac{r^n}{n!}
   \max_b|\psi_b^{(n)}(0)-\widetilde\psi_b^{(n)}(0)|
 \le \frac{(r/R)^3}{c(1-r/R)}
       \|\mathscr S(\psi)-\mathscr S(\widetilde\psi)\|_R.
\end{equation}
The same local real analytic inverse allows the gap and curvatures to
vary when the gap and both action Hessians are retained as finite
coordinates.  It concerns local contact germs and a strong norm on their
holomorphic extensions, not a weak-density or full-table statistical norm.
\end{theorem}

\begin{proof}
Fix the base point and abbreviate $L=L_{\psi^*}$, $A=A_{\psi^*}$ and
$K=K_{\psi^*}$.  We construct a bounded inverse for $L$, including its
lower-order couplings.  Choose one finite $m\ge4$ so large that
\begin{equation}\label{eq:v59-tail-choice}
                q_m:=\frac{2W}{a_*}\frac{\rho^m}{1-\rho^m}<1.
\end{equation}
On $\mathscr X_{R,m}$, Lemma~\ref{lem:v59-envelope-operator} gives
\begin{equation}\label{eq:v59-tail-inverse}
 L_{\rm tail}^{-1}
      =\sum_{j\ge0}(-A^{-1}K)^jA^{-1},\qquad
 \|L_{\rm tail}^{-1}\|\le\frac{a_*^{-1}}{1-q_m}.
\end{equation}
This norm-convergent series acts on actual functions.  It is not the
finite nilpotent formula with an unbounded number of unspecified factors.

Let $P_m$ be Taylor projection onto the pairs of polynomials of degrees
$3,\ldots,m-1$.  It is bounded on $\mathscr X_R$ by Cauchy's
coefficient estimate.  The induced finite map
\[
       L_{\rm low}=P_mL|_{P_m\mathscr X_R}
\]
is block lower triangular.  Its diagonal blocks are exactly
$M_3,\ldots,M_{m-1}$.  One way to see this is to differentiate the
actual finite-jet identity proved in
Lemma~\ref{lem:v27-smooth-jet-factorization} and
Proposition~\ref{prop:v22-last-jet-block}.  It also follows directly
from \eqref{eq:v59-full-envelope}: $w_{b,i}(0)=1$ and
$x_{b,i}'(0)=\ell_{b,i}$, so the new degree is the same weighted
geometric sum of visits.  Analytic continuation of this finite identity,
or complex linearity at the real base, gives the complex statement.
All these finitely many blocks are invertible.  Therefore $L_{\rm low}$
has a bounded inverse on its finite-dimensional space, without requiring
its strictly lower part to be small.

For an arbitrary $y\in\mathscr X_R$, define
\begin{equation}\label{eq:v59-full-linear-inverse}
 \eta_{\rm low}=L_{\rm low}^{-1}P_my,\qquad
 \eta_{\rm tail}=L_{\rm tail}^{-1}(y-L\eta_{\rm low}).
\end{equation}
The second argument lies in $\mathscr X_{R,m}$ because its projection
under $P_m$ vanishes.  These formulas solve $L\eta=y$ uniquely and
boundedly.  For example, if $B=\|L_{\rm low}^{-1}P_m\|$, then
\begin{equation}\label{eq:v59-complete-operator-bound}
 \|L^{-1}\|\le B+\frac{a_*^{-1}}{1-q_m}(1+\|L\|B).
\end{equation}
Only one fixed finite low-order problem has been inverted.  All remaining
orders are handled together by \eqref{eq:v59-tail-inverse}.

For completeness the nonlinear step can be made by contraction rather
than a formal inverse expansion.  Set
$\mathcal N(\eta)=\mathscr S(\psi^*+\eta)-\mathscr S(\psi^*)-L\eta$.
It is holomorphic and $D\mathcal N(0)=0$.  On a small ball arrange
$\|L^{-1}D\mathcal N\|<1/2$.  For $y$ small the equation
\[
                   \eta=L^{-1}y-L^{-1}\mathcal N(\eta)
\]
preserves this ball and is a contraction.  Its iterates converge locally
uniformly and holomorphically in $y$.  This gives the local holomorphic
inverse.  Shrinking a convex ball also gives
$\|D\mathscr S-L\|<(2\|L^{-1}\|)^{-1}$ there.  Integration on the
segment joining two points proves \eqref{eq:v59-nonlinear-bilip}, with
$c=(2\|L^{-1}\|)^{-1}$ and a local bound for the upper constant.
Real data give real inverses by uniqueness and conjugation symmetry.

Every tail $\mathscr X_{R,N+1}$ is invariant under $L$.
It is also invariant under $L^{-1}$: if the output has zero jets through
$N$, finite triangular induction with the invertible $M_n$ makes the
input jets through $N$ zero.  Thus $L$ and $L^{-1}$ descend to mutually
inverse quotient operators.  Their norms are bounded by the parent
operator norms, uniformly in $N$.  The same construction, or continuity
of the derivative and its inverse on a smaller parameter neighborhood,
gives local uniformity in the base pair.  This statement is about
analytic quotient norms, not maximum norms of unweighted derivatives.
Cauchy's estimate at zero and \eqref{eq:v59-nonlinear-bilip} then give
\eqref{eq:v59-all-coefficients} by summing the geometric series.

Finally, let the leading variables vary near their positive base values.
The protected-domain construction is holomorphic in $g,\kappa_0,\kappa_1$
as well.  Write $\psi_b=\kappa_bu^2/2+v_b$, with $v\in\mathscr X_R$,
and similarly subtract $a_bu^2/2$ from each action.  The leading map
$(g,\kappa_0,\kappa_1)\mapsto(g,a_0,a_1)$ has the local analytic
inverse already given by
\[
 \gamma=\operatorname{arsinh}(g\sqrt{a_0a_1}),\qquad
 c_b=\cosh\gamma\sqrt{a_b/a_{1-b}},\qquad \kappa_b=(c_b-1)/g,
\]
with the positive branches near the base.  In the finite-variable and
$\mathscr X_R$ splitting the derivative is block lower triangular:
the finite leading block and the infinite tail block are invertible.
The same local contraction proves the stated joint inverse.
\end{proof}

The theorem supplies a local analytic inverse for the entire action, not
just its diagonal response.  Its additional input is concrete: the norm
controls the holomorphic extensions on one fixed disc.  Smallness of a
real-interval $C^k$ error or total variation does not supply that norm.
The radius and the inverse constants may depend on the analytic base
pair, its gap, curvature margins and extension bounds.  The result does
not extend a local germ to an entire obstacle, select a discrete global
matching branch, or give a uniform acquisition cost.  Those remain the
separate geometric and statistical statements proved elsewhere.
Neither $K$ nor the full inverse is claimed to improve holomorphic radius;
only the finite-rank approximation and the tail bounds stated above are
used.  In particular the arbitrary triangular-matrix comparison does not
refute this result: an actual action derivative has the additional
multiplier-and-contracted-visits structure \eqref{eq:v59-full-envelope}.
